% SPDX-License-Identifier: Apache-2.0
\section{The decision, and what is missing from it}

An issuer authorising a card transaction has three actions, not two: approve, decline, or
challenge the cardholder through 3-D Secure. The middle action is what makes the problem
tractable at scale --- routing even one percent of authorisation volume to human analysts is
not operationally possible, while a step-up challenge is --- and it is the action that
published fraud-detection work almost never certifies. Models are reported by AUC or average
precision on a held-out fold. Neither is a statement about a decision, and neither carries a
finite-sample guarantee.

\textbf{What we contribute.} A three-valued rule whose \emph{abstention band} carries a
distribution-free, finite-sample risk certificate, evaluated under a pre-registration frozen
before any model was fitted; and an honest measurement of where two quantum approaches sit
against a tuned classical baseline on the same decision.

\textbf{Both quantum arms returned negative results.} A fidelity quantum kernel was rejected
by its own a-priori screens before it was run on the task. A matrix-product-state classifier
ran and did not beat gradient boosting. We report both in the body rather than the appendix,
because a screen that only ever passes is not a screen, and because the pre-registration
committed us to reporting whatever it returned.

\textbf{What we do not claim.} No quantum advantage of any kind. No portfolio-level PSD2
compliance figure --- the SCA-RTS ceilings govern exemption eligibility across a whole
portfolio, not a decline threshold on a fraud-enriched research benchmark, and treating them
as a threshold selector is a category error we made and withdrew. No claim that the
certificate binds the \emph{unconditional} false-decline rate: it does not, and
Section~\ref{sec:method} says why that distinction is the point rather than a caveat.
